
\subsection{Input Data}
The analysis was performed on the provided gene-level count matrix. Low-quality samples and conditions with an insufficient number of biological replicates were filtered out prior to analysis, and genes with very low expression across samples were removed to improve statistical power.

\subsection{Differential Gene Expression Analysis}

To ensure robust detection of differentially expressed genes (DEGs), a multi-method consensus approach was applied, combining complementary statistical frameworks on the normalised count data:

\begin{enumerate}
  \item \textbf{limma-voom} (Ritchie et al., 2015): models log-transformed data with precision weights derived from the mean--variance trend and employs empirical Bayes shrinkage of variance estimates.
  \item \textbf{edgeR} (Robinson et al., 2010): fits a negative binomial model and performs quasi-likelihood F-tests to account for overdispersion.
  \item \textbf{DESeq2} (Love et al., 2014): uses a negative binomial generalised linear model with shrinkage of dispersion parameters and log$_2$ fold-change estimates via \textit{apeglm} (Zhu et al., 2019).
\end{enumerate}

The p-values from all applicable methods were combined using Stouffer's Z-score meta-analysis approach, with p-values constrained to [10$^{-10}$, 0.95] to avoid numerical instabilities. Combined p-values were corrected for multiple testing using the Benjamini-Hochberg (BH) False Discovery Rate. Genes were classified as DEGs if: (i) combined adjusted p-value $< 0.05$, and (ii) absolute log$_2$ fold-change $\geq$ log$_2$(1.2).

\subsection{Functional Enrichment Analysis}

Functional enrichment analysis was performed using an over-representation analysis (ORA) approach against the local annotation database. Analyses were conducted separately for up-regulated and down-regulated DEG sets.

\textbf{Over-representation analysis (ORA)} tested whether DEGs were significantly over-represented in curated biological term sets using hypergeometric tests against multiple databases: Gene Ontology (GO --- Biological Process, Molecular Function, Cellular Component), KEGG pathways, Reactome pathways, WikiPathways, and MSigDB gene set collections (Hallmark, C2, C5, C6, C7).

Only results with Benjamini-Hochberg adjusted p-value $\leq$ 0.05 are reported. In dotplots, dot size reflects the gene ratio (proportion of DEGs annotated to that term) and dot colour represents statistical significance ($-\log_{10}$ adjusted p-value).
